\section{过渡态计算}
化学反应经历的起止点是反应物和产物的最稳定态。而两者之间衔接势能面\textrm{(Potential Energy Surface, PES)}的最高\textcolor{red}{鞍点}，称为过渡态\textrm{(Transition State)}。

\subsection{基本过渡态}

\subsubsection{反应路径与弹性带方法}
采用推拉弹性带方法\textrm{(Nudged Elastic Band, NEB)}方法，用4个影像计算反应路径，模拟1个\ch{Si}原子在15-原子的初级超晶胞表面的自扩散。

所有的化学反应都是由反应物经历过渡态到达产物。模拟的方法之一就是\text{NEB}\upcite{JCP113-9978_2000}。\textrm{NEB}在反应物和产物之间，插入若干影像点，模拟得到最小能量路径\textrm{(Minmum Energy Path, MEP)}，即反应路径。影像点的数目一般取为4-20个，总体上包括初态结构(反应物)到终态结构(产物)的变化过程。本算例中，选择4个影像，并选择稀疏基组。如果选择更多的影像，计算开销更大，收敛也会更困难。影像之间类似弹簧上的各环节，彼此关联，一直延伸到两头的初态结构和终态结构。这些彼此关联的影像形成了弹性的带。

带上的受力分解为两部分，传统的垂直于带$F_{i,\perpendicular}$的真实力和平行于带的弹性力$F_{i,\mathrm{spring}}$，称为\textrm{(Nudging)}。根据带上的受力$F_i$，对带的弛豫，即可得到反应路径。

本练习中，采用\textrm{NEB}模拟15-原子空位超晶胞如图所示，蓝色和海青色原子在晶胞中，红色表示空位。

在\textrm{NEB}计算中，蓝色位置的原子，如图所示，将移动到红色的空位。这就是金刚石型立方\ch{Si}原子晶胞中的自扩散，\upcite{PRB59-3969_1999}属于简单模拟。在练习中，通过自扩散研究反应路径。除了\textrm{NEB}方法，还有其他的过渡态方法。
\textrm{POSCAR}文件中，空位(\textrm{V1}和\textrm{V2})结构和影像都预先弛豫好了，因此只在部分\textrm{NEB}步中会需要弛豫处理。

通过\textrm{ASE}模块，构建\textrm{NEB}影像构建脚本。

由此产生的结构文件\textrm{POSCAR.V1}和\textrm{POSCAR.V2}都未经弛豫，在进行影像插值之前，必须先进行弛豫，相应的\textrm{INCAR}文件

弛豫结束后，编写\textrm{Python}脚本，应用\textrm{ASE}的插入影像功能，完成四个影像的插入。

得到包括初始构型、终了构型和插值影像构型在内的全部\textrm{POSCAR}文件

\textrm{00/POSCAR}文件

\textrm{01/POSCAR}文件

\textrm{02/POSCAR}文件

\textrm{03/POSCAR}文件

\textrm{04/POSCAR}文件

\textrm{05/POSCAR}文件

各结构所需的单点计算文件\textrm{INCAR.sp}

\textrm{NEB}计算文件\textrm{INCAR.neb}

$\vec k$空间网格剖分文件\textrm{KPOINTS}文件

\textrm{INCAR.sp}用于单点计算，而\textrm{INCAR.neb}用于\textrm{NEB}计算，两者的内容相同。

前四个计算控制参数用于电子步的最小化:~\texttt{ENCUT}确定平面波基组的数量；\texttt{PREC}确定计算对应的精度；\texttt{EDIFF}定义电子步自洽迭代的收敛标准，负值表示采用力收敛，单位是$\mathrm{eV\AA^{-1}}$，正值表示采用能量收敛。这里推荐力收敛。\textrm{Fermi}能积分展宽的计算控制参数是\texttt{ISMEAR}和\texttt{SIGMA}

其余计算控制参数用于离子弛豫设置。\texttt{NSW}是许可的最大离子步，对于单点能计算，该值为0；\texttt{EDIFFG}是离子弛豫收敛标准。当\texttt{IBRION}=1，结构弛豫采用准\textrm{Newton}算法。\texttt{POTIM}是离子弛豫的离子步长；\texttt{ISIF}=2只允许离子内坐标弛豫，但晶胞承诺书和晶格形状固定。计算过程中，可能会出现提示\textcolor{red}{\textbf ``BRIONS problem: POTIM should be increased''}，这表明收敛困难，建议增加\texttt{POTIM}的值或者改变\textrm{IBRION}的算法。

最后两个计算控制参数确定\textrm{NEB}，\texttt{IMAGES}设定采用\textrm{NEB}计算的影像(反应中间结构)的数目，本算例中数值为4\footnote{请注意:~并行计算时，\textrm{MPI}的数目必须影像的倍数。}\texttt{SPRING}=5是对\textrm{NEB}的弹性带设定弹簧倔强系数。

$\vec k$空间网格布点只有$\Gamma$点。

本算例中，有这些参数足以快速完成模拟计算，要获得更高质量的计算结果，可以增加$\vec k$空间网格剖分点和超晶胞的大小，直到最后的结果，如形成能收敛。离子步的收敛标准\texttt{EDIFFG}也要提升，如$-0.01\mathrm{eV\AA^{-1}}$。

在\textrm{NEB}计算中，共有六个计算目录，包括初始态(00)，依次四个影像态(01,02,03,04)和终了态(05)，在对应目录给出对应的\textrm{POSCAR}文件。\textrm{INCAR}文件中，影像数目用\texttt{IMAGES}来标记。

\textrm{NEB}计算所执行的\textrm{VASP}命令，在各态目录的上层父目录进行，然后\textrm{VASP}在各子目录执行相应的计算。完成准备工作后，执行\textrm{Bash}脚本，进行\textrm{VASP}计算。

直到\textrm{NEB}计算完成，获得了\ch{Si}原子扩散到临近空位的\textrm{NEB}影像的而能量。在讨论反应路径的时候，同样还需要起始态和终了态的结构信息，不过\textrm{NEB}计算中，起始态和终了态是按照静态处理的。这两个计算需要单独执行\textrm{Bash}环境下的\textrm{VASP}M命令。

为了节省计算开销，这里得到的能量与实验结果不匹配。如果要提升计算精度，有必要增加$\vec k$空间网格剖分密度，如$4\times4\times4$。推荐最终的单点计算采用四面体积分方案以得到好的计算精度，这里给出修改后的\textrm{INCAR}文件供参考。

计算结束，可以通过\textrm{Py4vasp}可视化原子结构的变化。

\ch{Si}原子大体上是沿着直线扩散的，考虑\ch{Si}超晶胞的交叉截面\textrm{(Cross-section)}，在反应通道上对应结构的变化，\ch{Si}原子由起始态位置向缺陷位置的扩散过程可以很容易地可视化。

以下是完成在二维平面上在反应通道标注影像结构的\textrm{Pyton}脚本。

对于这些影像，绝大多数原子留在原来的位置，扩散的\ch{Si}原子由起始位置向着空位移动的过程非常显眼。在此过程中，超晶胞结构也会发生变化，每个影像间隔差不多距离。采用\textrm{Py4vasp}可以看到反应通道中不同影像的能量变化。

得到的反应通道有一定的对称结构，峰位大约在$\sim0.7\mathrm{eV}$。如果要进一步研究过渡态，则需要在原有四个影像的基础上，继续增加影像，逼近完整的过渡态。需要说明的是，影像越少，计算的收敛性质会更好。而且如果选择太多影像，计算开销也很大。

通过这个\ch{Si}原子扩散到空位的例子，可以看到，对于过渡态的计算是复杂的。计算中有时候选择反应中间体影像作为寻找过渡态的试探结构。下面介绍如何优化试探结构。

\subsubsection{改进双粒子方法}
